% =============================================================================
\section{Introduction}
\label{sec:intro}
% =============================================================================
Bindings are essentially wrapper libraries that bridge two
programming languages, so that a software module developed in one
language can be used in code written in another language, exploiting
(i) the strengths of both languages, and (ii) the availability of
modules developed in the former language, if such exist. C++ Python
bindings enable the invocation of C++ functions from Python code
taking advantage of Python's quick development cycles and C++'s high
performance.
%
% =============================================================================
\subsection{C++ Python Bindings}
\label{ssec:intro:cgal}
% =============================================================================
Although Python is sufficiently efficient for many tasks, low-level
code written in Python tends to be too slow, largely due to Python's
extremely dynamic nature. In particular, low-level computational loops
are simply infeasible. In additiona, the magnitude of existing and
well-tested code in Fortract, C, and C++, is enourmous. Python
provides the means to acces its internals from these languages via the
Python/C API, which enables the exploitation of existing code in
Fortrant, C and C++, as well as the replacement of critical sections
written in these languages when speed is essential.

Wrapping existing code has traditionally been the domain of Python
experts due to the steep learning curve, which characterizes its Python/C
API. Although using such wrappers is possible without ever knowing
their internals, providing such wrappers creates a sharp line between
developers using Python and developers using C or C++ (while the
Python/C API). Several tools that ease the creation of such wrappers
have been proposed and introduced; A sample of such tools are listed In
Section~\ref{sec:comp}.

% =============================================================================
\subsection{\cmake}
\label{ssec:intro:cmake}
% =============================================================================
\cmake{}~\cite{mh-mc-08} is an open-source, cross-platform suite of tools designed to
build, test, and package software. The suite of \cmake{} tools were
created in response to the need for cross-platform build environments
for open-source projects. However, since its creation it has grown and
become powerful, rich, and flexible. It is now widely spread and used
by major commercial psoftware products for their build and test
environments.

\cmake{}, for the most part, is a cross-platform build system
generator.

It is used to control the software building process using
simple platform and compiler independent configuration files, and
generate native build environments of the user choice. For example,
once the \cmake{} configuration files are in place for a certain
project, \cmake{} can be used to generate \myLstinline{Makefile} files
required by the \myLstinline{make} utility of the GNU suite to compile
a project on a Unix system using \myLstinline{g++}.

Users build a
project by using \cmake{} to generate a build system for a native tool
on their platform.  The build process of a project is specified with
platform-independent \cmake{} text files, named
\myLstinline{CMakeLists.txt}, included in each directory of the source
tree of the project. Perhaps the most popular tool in the suite is the
\myLstinline{cmake} executable.  It is the command-line interface of
the cross-platform build system generator of \cmake. For graphical
user interfaces that may be used in place of cmake



To build a software project with CMake, Generate a Project Buildsystem. Optionally use cmake to Build a Project, Install a Project or just run the corresponding build tool (e.g. make) directly.

The other actions are meant for use by software developers writing scripts in the CMake language to support their builds.

we exploit this capability to conveniently generate the C++ Python bindings.

For command-line interfaces to the CMake testing and packaging facilities, see ctest and cpack.

% =============================================================================
\subsection{\cgal}
\label{ssec:intro:cgal}
% =============================================================================
\cgal{} (Computational Geometry Algorithm Library) is an open source
software project that provides easy access to efficient and reliable
implementations of geometric algorithms in the form of a C++
library~\cite{cgal:eb-21}. \cgal{} has evolved through the years and
now represents the state of the art in applied computational geometry
software. The software modules of \cgal{} rigorously adhere to the
generic programming paradigm to overcome problems encountered when
effective computational geometry software is implemented; they are
stable despite the use of (inexact) floating point arithmetic, found
in standard hardware, and they are complete as they handle all
degenerate cases, which typically are in abundance.

When running Python code that use bindings for some C++ modules, one
or more libraries that provide the bindings must be accessible.  A
software module that adheres to the generic programming paradigm
consists of function and class templates; these templates are
instantiated at compile time of the binding libraries. In other words,
the type of the C++ objects that are bound, that is, instances
(instantiated types) of C++ function and class templates, must be
known when the bindings are generated.

Generic programming identifies an abstraction that consists of a
formal hierarchy of polymorphic abstract requirements on data types,
referred to as \emph{concepts}, and a set of classes that conform
precisely to the specified requirements, referred to as
\emph{models}~\cite{a-gps-99}. When a class template is instantiated,
each one of its template parameters is substituted with a model of a
concept associated with the template parameter.

Instantiated types in \cgal{} are characterized by long instantiation
chains of C++ class templates. An instantiated type is a template that
has one or more template parameters and every parameter is substituted
by another type that is typically an instance of another template.
%
While the number of models of most concepts is small, the number of
potential types of objects that could be bound, and thus must be
supported by the bindings, is enormous. Offering bindings for all
these types in advance is practiaclly impossible. Moreover, in some
cases models need to be extended with types provided by the
user. Naturally, bindings for user-extended models must be generated
dynamically.

\cgal{} uses \cmake{} for the build process.  Since version~5.0,
\cgal{} is header-only be default, which means that there is no need
to compile any library before \cgal{} can be used by an
application. However, some dependencies of \cgal{} might need to be
installed, and the configuration of \cgal{} must be carried out, which
consists of a single application of xxx cmake.

% =============================================================================
\subsection{Results}
\label{ssec:intro:results}
% =============================================================================
We introduce C++ Python bindings that enables the convenient,
efficient, and reliable use of \cgal{} software modules in Python's
code. There are different tools that facilitate the creation of C++
Python bindings. We present a short study that compares three main
tools, which leads to the method of choice.
%
We present a system that dynamically generates bindings of desired
objects according to user prescriptions, which enables the convenient
use of any subset of bound object types concurrently.  With our system
it is possible to generate, for example, a single library that
contains bindings for instances of different \cgal{} class templates,
e.g., an instance of the 2D arrangement class template and an instance
of the 2D triangulation class template. If several instances of the
same class template must be bound, several corresponding libraries are
generated, one for each instance.
%
Finally, the binding themselves are implemented in C++ and their
implementation exploits advanced features in C++. This results with
elegant and compact code that is easy to maintain and extend with
additional bindings.

% =============================================================================
\subsection{Outline}
\label{ssec:intro:outline}
% =============================================================================
The rest of this paper is organized as follows. A study that compares
three main methods for implementing C++ Python bindings is presented
in Section~\ref{sec:comp}. Description of the bindings generated by
our system and instructons for dynamically generating bindings are
presented in Section~\ref{sec:bind}.
